\documentclass{tufte-book}

\hypersetup{colorlinks}% uncomment this line if you prefer colored hyperlinks (e.g., for onscreen viewing)

\usepackage[superscript,biblabel]{cite}

\geometry{a5paper,left=20mm,top=15mm,headsep=1\baselineskip,textwidth=95mm,textheight=180mm,headheight=\baselineskip,marginparwidth=20mm}

%%
% Book metadata
\title{A Tufte-Style Book\thanks{Thanks to the Tufte-LaTeX\ Developers}}
\author{Susan E. McGregor - susanemcg}
\publisher{Tow Center for Digital Journalism}


% use fontspec package and set titles to Georgia
\usepackage{fontspec}
\newfontfamily\titlefont{Georgia}

% add color package and define appropriate blues for title page
\usepackage{color}
\definecolor{darkblue}{RGB}{17,94,170}
\definecolor{lightblue}{RGB}{188,229,245}

% framed helps with graphics
\usepackage{framed}

% suppresses names and figure numbers on graphics
\usepackage{etoolbox}
\makeatletter
\patchcmd{\@caption}
  {\noindent\csname fnum@#1\endcsname: \ignorespaces}
  {}
  {}{}
\makeatother

%%
% Just some sample text
\usepackage{lipsum}

%%
% For nicely typeset tabular material
\usepackage{booktabs}

%%
% For graphics / images
\usepackage{graphicx}
\setkeys{Gin}{width=\linewidth,totalheight=\textheight,keepaspectratio}
\graphicspath{{graphics/}}

% The fancyvrb package lets us customize the formatting of verbatim
% environments.  We use a slightly smaller font.
\usepackage{fancyvrb}
\fvset{fontsize=\normalsize}

%%
% Prints argument within hanging parentheses (i.e., parentheses that take
% up no horizontal space).  Useful in tabular environments.
\newcommand{\hangp}[1]{\makebox[0pt][r]{(}#1\makebox[0pt][l]{)}}

%%
% Prints an asterisk that takes up no horizontal space.
% Useful in tabular environments.
\newcommand{\hangstar}{\makebox[0pt][l]{*}}

%%
% Prints a trailing space in a smart way.
\usepackage{xspace}


% Prints the month name (e.g., January) and the year (e.g., 2008)
\newcommand{\monthyear}{%
  \ifcase\month\or January\or February\or March\or April\or May\or June\or
  July\or August\or September\or October\or November\or
  December\fi\space\number\year
}


% Prints an epigraph and speaker in sans serif, all-caps type.
\newcommand{\openepigraph}[2]{%
  %\sffamily\fontsize{14}{16}\selectfont
  \begin{fullwidth}
  \sffamily\large
  \begin{doublespace}
  \noindent\allcaps{#1}\\% epigraph
  \noindent\allcaps{#2}% author
  \end{doublespace}
  \end{fullwidth}
}

% Inserts a blank page
\newcommand{\blankpage}{\newpage\hbox{}\thispagestyle{empty}\newpage}

\usepackage{units}

% Typesets the font size, leading, and measure in the form of 10/12x26 pc.
\newcommand{\measure}[3]{#1/#2$\times$\unit[#3]{pc}}

% Macros for typesetting the documentation
\newcommand{\hlred}[1]{\textcolor{Maroon}{#1}}% prints in red
\newcommand{\hangleft}[1]{\makebox[0pt][r]{#1}}
\newcommand{\hairsp}{\hspace{1pt}}% hair space
\newcommand{\hquad}{\hskip0.5em\relax}% half quad space
\newcommand{\TODO}{\textcolor{red}{\bf TODO!}\xspace}
\newcommand{\ie}{\textit{i.\hairsp{}e.}\xspace}
\newcommand{\eg}{\textit{e.\hairsp{}g.}\xspace}
\newcommand{\na}{\quad--}% used in tables for N/A cells
\providecommand{\XeLaTeX}{X\lower.5ex\hbox{\kern-0.15em\reflectbox{E}}\kern-0.1em\LaTeX}
\newcommand{\tXeLaTeX}{\XeLaTeX\index{XeLaTeX@\protect\XeLaTeX}}
% \index{\texttt{\textbackslash xyz}@\hangleft{\texttt{\textbackslash}}\texttt{xyz}}
\newcommand{\tuftebs}{\symbol{'134}}% a backslash in tt type in OT1/T1
\newcommand{\doccmdnoindex}[2][]{\texttt{\tuftebs#2}}% command name -- adds backslash automatically (and doesn't add cmd to the index)
\newcommand{\doccmddef}[2][]{%
  \hlred{\texttt{\tuftebs#2}}\label{cmd:#2}%
  \ifthenelse{\isempty{#1}}%
    {% add the command to the index
      \index{#2 command@\protect\hangleft{\texttt{\tuftebs}}\texttt{#2}}% command name
    }%
    {% add the command and package to the index
      \index{#2 command@\protect\hangleft{\texttt{\tuftebs}}\texttt{#2} (\texttt{#1} package)}% command name
      \index{#1 package@\texttt{#1} package}\index{packages!#1@\texttt{#1}}% package name
    }%
}% command name -- adds backslash automatically
\newcommand{\doccmd}[2][]{%
  \texttt{\tuftebs#2}%
  \ifthenelse{\isempty{#1}}%
    {% add the command to the index
      \index{#2 command@\protect\hangleft{\texttt{\tuftebs}}\texttt{#2}}% command name
    }%
    {% add the command and package to the index
      \index{#2 command@\protect\hangleft{\texttt{\tuftebs}}\texttt{#2} (\texttt{#1} package)}% command name
      \index{#1 package@\texttt{#1} package}\index{packages!#1@\texttt{#1}}% package name
    }%
}% command name -- adds backslash automatically
\newcommand{\docopt}[1]{\ensuremath{\langle}\textrm{\textit{#1}}\ensuremath{\rangle}}% optional command argument
\newcommand{\docarg}[1]{\textrm{\textit{#1}}}% (required) command argument
\newenvironment{docspec}{\begin{quotation}\ttfamily\parskip0pt\parindent0pt\ignorespaces}{\end{quotation}}% command specification environment
\newcommand{\docenv}[1]{\texttt{#1}\index{#1 environment@\texttt{#1} environment}\index{environments!#1@\texttt{#1}}}% environment name
\newcommand{\docenvdef}[1]{\hlred{\texttt{#1}}\label{env:#1}\index{#1 environment@\texttt{#1} environment}\index{environments!#1@\texttt{#1}}}% environment name
\newcommand{\docpkg}[1]{\texttt{#1}\index{#1 package@\texttt{#1} package}\index{packages!#1@\texttt{#1}}}% package name
\newcommand{\doccls}[1]{\texttt{#1}}% document class name
\newcommand{\docclsopt}[1]{\texttt{#1}\index{#1 class option@\texttt{#1} class option}\index{class options!#1@\texttt{#1}}}% document class option name
\newcommand{\docclsoptdef}[1]{\hlred{\texttt{#1}}\label{clsopt:#1}\index{#1 class option@\texttt{#1} class option}\index{class options!#1@\texttt{#1}}}% document class option name defined
\newcommand{\docmsg}[2]{\bigskip\begin{fullwidth}\noindent\ttfamily#1\end{fullwidth}\medskip\par\noindent#2}
\newcommand{\docfilehook}[2]{\texttt{#1}\index{file hooks!#2}\index{#1@\texttt{#1}}}
\newcommand{\doccounter}[1]{\texttt{#1}\index{#1 counter@\texttt{#1} counter}}

% Generates the index
\usepackage{makeidx}
\makeindex

\begin{document}

% Front matter
\frontmatter

% r.1 blank page
\blankpage


% r.3 full title page
%\maketitle
%Using minipages to mimic existing layout for title page

\begin{minipage}[t][1\textheight]{0.40\textwidth}
\vspace{0pt}
\begin{flushleft}
\small\textsf{Tow Center for Digital Journalism\\
A Tow/Knight Report}
\end{flushleft}
\vfill
\includegraphics[width=.4\textwidth]{graphics/ColumbiaLogo.pdf}
\end{minipage}%
\hfill
\begin{minipage}[t][1\textheight]{70mm}
\vspace{0pt}
\Huge\textcolor{darkblue}{\textbf{{\titlefont DIGITAL\\ SECURITY\\
AND\\
SOURCE\\ PROTECTION\\
FOR\\
JOURNALISTS\\
\textcolor{lightblue}{
SUSAN\\
MCGREGOR\\
}
}
}}
\vfill
\small\textsf{Funded by the Tow Foundation\\
and the John S. and James L. Knight Foundation}
\end{minipage}%


% v.4 copyright page
\newpage
\null
\vfill
\begin{fullwidth}
\textsf{\textbf{Acknowledgements}} 
\\[1.0cm]
\noindent\textit{This paper is the product of countless conversations, encounters, recommendations and insights provided to me over the past year by members of the digital security and journalism communities.}\\[0.5cm]

\noindent\textit{Though it would be impossible to name everyone whose insight and expertise has informed this work, my deepest thanks go to Franzi Roesner, Carey Shenkman, Nabiha Syed, Seamus Tuohy, Chris Walker and Carol Waters for providing invaluable feedback and direction on early versions of this paper; your expertise and insight was invaluable. Also to Oktavia Jonsdottir, Scott Klein, Quinn Norton, Jeremy Singer-Vine and Jennifer Valentino-DeVries, who shared their ideas and work processes with me. And a very special thanks to Peter G. Neumann, who offered me the long view on digital security and directed me to essential resources and experts on the subject.} \\[0.5cm]

\noindent\textit{This project would also not have been possible without the financial and intellectual support of my colleagues at Columbia Journalism School, and the Tow Center for Digital Journalism--particularly Emily Bell and Taylor Owen. Likewise, I am very grateful to both OpenITP and TacticalTech--organizations whose generosity in hosting me in the fall of 2013 made possible so many enlightening interactions with the digital security community. Without the openness of these organizations and their associates, this work would not have been possible.} \\[0.5cm]

\noindent\textit{And finally, many grateful thanks to my talented, enthusiastic, and insightful illustrator and collaborator, Matteo Farinella, for truly making the invisible, visible.} \\[1.0cm]

\noindent \small{Copyright \copyright\ \the\year\ \thanklessauthor}

\noindent \small{This volume produced with the assistance of the \textsf{tufte-\LaTeX{}} class.}

\noindent\textit{First printing, June 2014}
%\monthyear}
\end{fullwidth}

% r.5 contents
\tableofcontents


% r.7 dedication
\cleardoublepage
~\vfill
\begin{doublespace}
\noindent\fontsize{16}{18}\selectfont\itshape
\nohyphenation
This paper is dedicated to all members of the digital security and journalism communities, whose shared commitment to transparency, skepticism towards power, and dedication to informing and empowering individuals everywhere has been an inspiration. I hope that this will be only the first of many collaborations between our equally scrappy, resilient, and adaptive professions. 
\end{doublespace}
\vfill
\vfill


% r.9 introduction
\cleardoublepage
\chapter*{Introduction}

Add introduction here.


%%
% Start the main matter (normal chapters)
\mainmatter

\chapter{A chapter title}
 Lorem ipsum dolor sit amet, consectetur adipiscing elit. Duis facilisis faucibus libero in tempor. Donec diam eros, elementum ac porta vitae, porttitor sit amet orci. Integer vitae ex at magna posuere volutpat at vitae dui. Aliquam efficitur ante ligula, at facilisis risus gravida ut. Etiam justo ligula, imperdiet in lectus sit amet, ultrices suscipit sapien. Nunc scelerisque semper dignissim. Curabitur eu accumsan erat. Ut nec ultricies ante. Nam nec vestibulum ligula.

Donec aliquet egestas elit elementum gravida. Cum sociis natoque penatibus et magnis dis parturient montes, nascetur ridiculus mus. Cras vitae augue ullamcorper, sodales lacus quis, sollicitudin erat. Sed laoreet rhoncus nibh, quis faucibus nibh feugiat vitae. Sed mollis, ante non ultrices imperdiet, metus metus iaculis risus, vitae blandit quam odio eu mauris. Curabitur a fermentum turpis. Aliquam sed odio quis augue posuere efficitur vitae sed augue\cite{Tufte2006}.

In suscipit, sapien vel varius pharetra, arcu turpis varius mi, sed luctus libero felis id lorem. Sed a sapien magna. Maecenas nunc quam, euismod at tincidunt eget, dignissim non ante. Phasellus condimentum dignissim ligula eu egestas. Ut faucibus laoreet pharetra. Nunc feugiat quam in ligula gravida aliquet. Cum sociis natoque penatibus et magnis dis parturient montes, nascetur ridiculus mus. Duis sem nisl, condimentum at dapibus non, sagittis eu velit. Ut porta turpis sem, elementum fringilla justo accumsan eget. Sed fringilla tempus quam a cursus. Fusce faucibus dignissim nunc eu lacinia. 


%%
% The back matter contains appendices, bibliographies, indices, glossaries, etc.



\backmatter

\bibliography{sample}
\bibliographystyle{unsrt}



\printindex

\end{document}

